\chapter{Implementación evolutiva y pruebas automatizadas}
\label{chap:implementacion}

\begin{objetivos}
Al finalizar este capítulo, el estudiante podrá construir un incremento vertical pequeño, separar pruebas unitarias, de integración y de aceptación, medir cobertura sin confundirla con calidad y definir puertas automáticas de validación.
\end{objetivos}

\section{Construir por incrementos verificables}

Una arquitectura solo existe de manera útil cuando sus decisiones pueden verse en código ejecutable. Un \emph{incremento vertical} atraviesa interfaz, lógica de aplicación y persistencia para entregar una capacidad observable. Es preferible implementar una reserva completa y pequeña que crear, por separado, muchas capas aún inconexas.

El ciclo recomendado es: formular una hipótesis, escribir una prueba que represente el comportamiento, implementar lo mínimo, refactorizar y registrar la evidencia. La prueba no sustituye el razonamiento de diseño; lo vuelve repetible.

\begin{figure}[htbp]
\centering
\begin{tikzpicture}[node distance=1.4cm and 1.0cm, every node/.style={font=\small}]
\node[decision] (h) {Hipótesis de comportamiento};
\node[activity, right=of h] (p) {Prueba que falla};
\node[activity, right=of p] (i) {Implementación mínima};
\node[activity, below=of i] (r) {Refactorización};
\node[evidence, left=of r] (e) {Evidencia reproducible};
\draw[arrow] (h)--(p); \draw[arrow] (p)--(i); \draw[arrow] (i)--(r);
\draw[arrow] (r)--(e); \draw[arrow] (e.west) -| (h.south);
\end{tikzpicture}
\caption{Ciclo corto de implementación y validación.}
\end{figure}

\section{Una estrategia de pruebas por riesgo}

La pirámide de pruebas es una guía económica, no una cuota fija. Las pruebas unitarias aíslan reglas rápidas; las de integración verifican contratos con base de datos, cola o sistema de archivos; las de extremo a extremo recorren el sistema desplegado. Cuanto más amplia sea la prueba, mayor suele ser su costo de preparación, tiempo y diagnóstico.

\begin{table}[htbp]
\centering
\caption{Selección de pruebas según el riesgo.}
\begin{tabularx}{\textwidth}{p{2.7cm}X X}
\toprule
\textbf{Nivel} & \textbf{Pregunta principal} & \textbf{Ejemplo en LibreReserva}\\
\midrule
Unitaria & ¿La regla produce el resultado correcto? & Rechazar un intervalo que termina antes de comenzar.\\
Integración & ¿La colaboración real cumple el contrato? & Guardar una reserva y detectar solapamiento en PostgreSQL.\\
Contrato & ¿Consumidor y proveedor interpretan igual la interfaz? & Verificar esquema y códigos de la API.\\
Aceptación & ¿El flujo entrega valor desde fuera? & Crear, consultar y cancelar una reserva por HTTP.\\
Rendimiento & ¿Se mantiene el objetivo bajo carga? & Percentil 95 menor al umbral con concurrencia definida.\\
Seguridad & ¿Los controles resisten usos adversos? & Negar acceso a reservas de otro usuario.\\
\bottomrule
\end{tabularx}
\end{table}

\section{Dobles, propiedades y datos de prueba}

Un doble de prueba reemplaza una dependencia para controlar el entorno. Un \emph{stub} responde valores preparados, un \emph{fake} ofrece una implementación simplificada y un \emph{mock} verifica interacciones. Deben usarse en los bordes; simular cada objeto interno acopla las pruebas a la implementación.

Las pruebas basadas en propiedades expresan invariantes para muchos datos generados. Por ejemplo: toda reserva aceptada tiene duración positiva y ninguna secuencia de creación puede dejar dos reservas activas solapadas para el mismo recurso. En Python, Hypothesis es una opción libre para generar casos y reducir automáticamente un fallo al ejemplo mínimo.

\begin{instalacion}
Se requiere Python 3.11 o superior. En un entorno virtual instale las herramientas libres:
\begin{lstlisting}[language=bash]
python3 -m venv .venv
source .venv/bin/activate          # Windows PowerShell: .venv\Scripts\Activate.ps1
python -m pip install fastapi uvicorn pytest pytest-cov httpx hypothesis
\end{lstlisting}
No use \texttt{sudo pip}. El entorno virtual hace reproducible el laboratorio y evita modificar el Python del sistema.
\end{instalacion}

\section{Ejemplo: contrato HTTP comprobable}

\begin{lstlisting}[language=Python,caption={API mínima con validación explícita.}]
from datetime import datetime
from fastapi import FastAPI, HTTPException, status
from pydantic import BaseModel, model_validator

app = FastAPI(title="LibreReserva")
reservas: list[dict] = []

class ReservaEntrada(BaseModel):
    recurso: str
    inicio: datetime
    fin: datetime

    @model_validator(mode="after")
    def intervalo_valido(self):
        if self.fin <= self.inicio:
            raise ValueError("fin debe ser posterior a inicio")
        return self

@app.post("/reservas", status_code=status.HTTP_201_CREATED)
def crear(entrada: ReservaEntrada):
    hay_cruce = any(
        r["recurso"] == entrada.recurso
        and entrada.inicio < r["fin"] and r["inicio"] < entrada.fin
        for r in reservas
    )
    if hay_cruce:
        raise HTTPException(status_code=409, detail="intervalo ocupado")
    nueva = {"id": len(reservas) + 1, **entrada.model_dump()}
    reservas.append(nueva)
    return nueva
\end{lstlisting}

\begin{lstlisting}[language=Python,caption={Pruebas externas del comportamiento HTTP.}]
from fastapi.testclient import TestClient
from app import app, reservas

cliente = TestClient(app)

def setup_function():
    reservas.clear()

def test_crea_y_rechaza_solapamiento():
    datos = {"recurso": "A-101",
             "inicio": "2026-08-10T08:00:00Z",
             "fin": "2026-08-10T09:00:00Z"}
    assert cliente.post("/reservas", json=datos).status_code == 201
    assert cliente.post("/reservas", json=datos).status_code == 409

def test_rechaza_intervalo_invertido():
    datos = {"recurso": "A-101",
             "inicio": "2026-08-10T10:00:00Z",
             "fin": "2026-08-10T09:00:00Z"}
    assert cliente.post("/reservas", json=datos).status_code == 422
\end{lstlisting}

\section{Cobertura, mutación y calidad}

La cobertura indica qué código fue ejecutado, no si fue correctamente comprobado. Una cobertura alta puede coexistir con aserciones débiles. La prueba de mutación cambia deliberadamente operadores o constantes: si la suite no detecta el cambio, falta sensibilidad. Para un curso corto es suficiente combinar cobertura de ramas, revisión de aserciones y una mutación manual de una regla crítica.

\begin{validacion}
Ejecute \texttt{pytest --cov=. --cov-branch --cov-report=term-missing}. Deben aprobarse todas las pruebas y cubrirse ambas decisiones del solapamiento. Cambie temporalmente \texttt{entrada.inicio < r["fin"]} por \texttt{<=}; documente qué prueba falla. Restaure el código antes de confirmar cambios.
\end{validacion}

\begin{ejercicio}
\textbf{Laboratorio 7.1 --- incremento vertical.} En Git y FastAPI implemente consulta y cancelación. Añada pruebas unitarias de dominio, integración HTTP y al menos una propiedad con Hypothesis. Entregue código, informe de cobertura y una nota sobre un defecto descubierto por las pruebas.

\textbf{Laboratorio 7.2 --- presupuesto de rendimiento.} Instale Locust con \texttt{python -m pip install locust}. Modele una carga reproducible, mida tasa de solicitudes, errores y p95, y grafique el resultado. Declare máquina, versión, datos y duración: sin esas condiciones la cifra no es comparable.
\end{ejercicio}

\section{Preguntas de revisión}

\begin{enumerate}
  \item ¿Qué riesgo quedaría invisible si solo se escribieran pruebas unitarias?
  \item ¿Por qué un porcentaje de cobertura no debe convertirse por sí solo en objetivo de aprendizaje?
  \item ¿Qué evidencia distingue una prueba inestable de un defecto intermitente real?
\end{enumerate}
